\setcounter{page}{1}
\phantomsection
\addcontentsline{toc}{section}{Abstract}
%ca 10 sätze
\begin{abstract}
 Die vorliegende Arbeit zeigt den methodischen Transfer des Frameworks der Systemic Functional Linguistics nach M.A.K. Halliday auf diachrone Korpuslinguistik. Ziel ist es, zu untersuchen, inwiefern durch das aus diesem Framework stammende \textit{transitivity system} zur Untersuchung verbaler Semantik auf korpusbasierte Rechtslinguistik übertragen werden kann. Die Datengrundlage der Untersuchung stellt dabei das mulitlayered Korpus \textit{konkorp} dar \citep{Thurau2025}, welches sechs deutsche Verfassungstexte enthält.
 Der Arbeit liegt folgende Forschungsfrage zugrunde:

 \begin{quote}
  Ist es möglich, verfassungsrechtliche Paradigmenwechsel durch eine systemisch-funktionale Analyse von Verbphrasen in Verfassungstexten zu identifizieren? Falls ja, wie lassen sich diese grammatisch klassifizieren?
 \end{quote}


 Durch die Zuordnung der im Korpus enthaltenen Verbgefüge zu einer von sechs funktionalen Kategorien von \textit{processes} aus der SFL, sowie die Untersuchung der \textit{circumstances} und \textit{participants} und der Veränderung im Zeitraum 1848--1968 wird dargestellt, ob und inwiefern ein Sprachwandel in Bezug auf Verben, demnach in Bezug auf die Sprache im deutschen Verfassungsrecht und somt letztlich im deutschen Verfassungsrecht stattgefunden hat.

 Die Arbeit enthält eine quantitativen Teil in Form einer Frequenzanalyse, die die Häufigkeit der Verbgefüge und deren Wandel darstellt. Des Weiteren wird in einem qualitativen Teil durch eine systemisch-funktionale Analyse ausgewertet, wie sich Verbsemantik als \textit{transitivity system} in Verfassungstexten darstellen lässt.
\end{abstract}
